\section{The analytic half: from prime sums on windows to zero-free half-planes}\label{sec:analytic}

This section follows the left column of Figure~\ref{fig:top}. Its result is that
Statement~\ref{st:every} follows from three inputs for each primitive character $\psi$: a
Landau-type criterion and a conditional square-root bound, both proved in the tree for every
``level'', and the existence of suitable families of auxiliary data on windows, which the
automorphic half supplies (\S\ref{sec:window}, \S\ref{sec:automorphic}).

\subsection{Primitive characters}

The analytic half works with primitive characters, packaged as a structure, and with the
proposition that all of their $L$-functions are zero-free on $\operatorname{Re}s>\frac12$:
\leanquote{Formalism.HC.PChar}
\leanquote{Formalism.HC.PChar.ZeroFreeAbove}
\leanquote{Formalism.HC.ZeroFreeAll}
\noindent Here \lean{Formalism.HC.ZeroFreeOn} \texttt{Q} requires \texttt{ψ.ZeroFreeAbove (1/2)} for
every \texttt{ψ} with \texttt{Q ψ} \stl{src/Formalism/SharpRestriction/RamificationFrobenius/TotientGaloisClasses/ExplicitFormulaErrorOfHostConditions.lean}{14}.
The passage back to all characters is the Euler-factor argument of \S\ref{sec:top} once more:
\leanquote{UpperTheory.Top.imprimitive}

\begin{note}
\lean{Formalism.HC.ZeroFreeAll} includes the trivial character modulo $1$ at $s=1$, where Mathlib's
$L$-function is \mathlib{riemannZeta} evaluated at its pole, a conventional finite value. The
statement therefore contains the assertion that this value is non-zero, which holds for Mathlib's
value (\S\ref{sec:top}). Statement~\ref{st:every}
excludes the pair ($\chi=1$, $s=1$) explicitly. \sti
\end{note}

\subsection{The prime sum}

The central analytic quantity is a weighted von Mangoldt sum with its expected main term removed:
\leanquote{Formalism.Sz}
\leanquote{Formalism.Gz}
\noindent That is,
\[
  S(x)=\sum_{n\le x}\Lambda(n)\Bigl(\frac{x}{n}-1\Bigr),\qquad
  G(x)=S(x)-x\log x+(1+\gamma)x .
\]
This is exactly the expression of remark item 5 \sta{5}. Its meaning is classical. Writing
$S(x)=\int_1^x\sum_{n\le t}\Lambda(n)n^{-1}\,dt$ and applying Perron's formula with the kernel
$x^{w}/(w(w-1))$ gives
\[
  S(x)=\frac1{2\pi i}\int_{(c)}\Bigl(-\frac{\zeta'}{\zeta}(w)\Bigr)\frac{x^{w}}{w(w-1)}\,dw\qquad(c>1).
\]
The double pole at $w=1$ contributes $x\log x-(1+\gamma)x$, which is precisely the term removed in
$G$. At $w=0$ the residue is $\log(2\pi)$. The non-trivial and trivial zeros give, for $x>1$,
\[
  G(x)=\log(2\pi)-\sum_{\rho}\frac{x^{\rho}}{\rho(\rho-1)}-\sum_{k\ge1}\frac{x^{-2k}}{2k(2k+1)},
\]
with non-trivial zeros counted with multiplicity. The zero sum converges absolutely and the
trivial-zero sum is $O(x^{-2})$ \stcs~\cite[Ch.~17]{davenport}\cite[Ch.~5]{iwaniec-kowalski}.
So $G(x)=O(x^{1/2})$ under the Riemann hypothesis, while a zero with real part $\beta>\frac12$
forces oscillations of size $x^{\beta}$. The tree uses twisted versions as well:
\leanquote{Formalism.HC.An.gChi}
\noindent For a non-principal character and $y\ge1$, the two terms are the weighted sum
$\sum_{n\le y}\chi(n)\Lambda(n)(y/n-1)$ and $y\,(L'/L)(1,\chi)$
\stl{src/Formalism/SharpRestriction/RamificationFrobenius/TotientGaloisClasses/Analytic.lean}{18}
\stl{src/Formalism/SharpRestriction/RamificationFrobenius/TotientGaloisClasses/Analytic.lean}{11}.
The latter cancels the residue contribution at $w=1$ \sti. A \emph{level} \lean{Formalism.HC.Level}
is a finite set \texttt{Λ.KEPT} of primitive characters with a natural number \texttt{o} and an
optional distinguished character \texttt{psK}. Its prime sum adds to $G$ the weighted real parts of
these twisted sums and their corrections, over the non-principal kept characters:
\leanquote{Formalism.HC.Level.Gm}
\noindent with weight $1$ for the distinguished character and $(3o+1)^{-1}$ for the others
\stl{src/Formalism/SharpRestriction/RamificationFrobenius/TotientGaloisClasses/ExplicitFormulaErrorOfHostConditions.lean}{36}.
The level used for $\psi$ is \lean{Formalism.HC.Level.cyc} $\psi$, built from the cyclic group
generated by $\psi$ \stl{src/Formalism/SharpRestriction/RamificationFrobenius/TotientGaloisClasses/ExplicitFormulaErrorOfHostConditions.lean}{121}.

\subsection{The Landau-type criterion on windows}

A set of reals \emph{has windows at scale} $\lambda$ if it contains intervals $[Y,Y^{\lambda}]$ with
$Y$ arbitrarily large:
\leanquote{Formalism.HC.HasWindows}
\noindent The criterion asserts that a one-sided bound $\Lambda.\mathrm{Gm}(x)\le x^{\theta}$ on such
windows excludes zeros to the right of $\theta$:
\leanquote{Formalism.HC.Level.FormBAt}
\noindent It is proved for every level and every $\lambda>1$:
\leanquote{UpperTheory.Comparison₂.formB}
\noindent The proof transports the level to a weighted family of characters with the same prime sum
and the same zero-freeness (\lean{UpperTheory.Comparison₂.level_family},
\leanloc{src/Development/AutomorphicArithmetic/PrimitiveCosets/LevelFamily.lean}{41}) and applies
\leanquote{Deduction.zeroFree_of_rpow_window}

\begin{statement}[Landau-type criterion]\label{st:landau}
Let $F$ be a finite family of positively weighted Dirichlet characters in which any principal
character has modulus $1$ (a level becomes one through
\lean{UpperTheory.Comparison₂.level_family}). Fix $0<\theta<1$ and $\delta>0$. Suppose there is a
real $C$ such that for every $Y_0$ there is $Y\ge Y_0$ with $\mathrm{Gm}_F(y)\le C\,y^{\theta}$ for
all $y\in[Y,Y^{1+\delta}]$. Then each $L$-function in $F$ is non-zero for
$\operatorname{Re}s>\theta$, in Mathlib's total-function convention.
\stl{src/Development/AutomorphicArithmetic/DedekindLFunctions/LandauTheorems/ZeroFreeOfRpowWindow.lean}{11}
\end{statement}

Two features distinguish this from the textbook statement. First, the bound is one-sided. In the
classical theory a one-sided bound excludes zeros through positivity: by Landau's theorem, a
Dirichlet series or Mellin transform of a function of constant sign is analytic to the right of a
real singularity. Applied to $C x^{\theta}-G(x)$, a one-sided bound on $G$ pushes every pole coming
from a zero $\rho$ to $\operatorname{Re}\rho\le\theta$, provided no real singularity intervenes
\stcs~\cite[Ch.~1 and~15]{montgomery-vaughan}. Second, the bound is required only on windows
$[Y,Y^{1+\delta}]$, not for all large $y$. We have not read the proof of
Statement~\ref{st:landau} (the lemma \lean{Deduction.lemmaL} that it calls) and do not know which
positivity it uses; this is question~Q4 in the list for outside experts.

\subsection{The conditional square-root bound}

The bound fed into Statement~\ref{st:landau} comes from auxiliary \emph{families}. A family
\texttt{Fm} (of type \lean{Formalism.BasicNotions.Fam} $\varepsilon$, or its variant
\lean{Formalism.MainNotions.FamP}) carries a set of windows \texttt{Win} and, for each $x$,
an index set of local and global data. The square-root bound is stated relative to a predicate
\texttt{Tied} on families:
\leanquote{Formalism.HC.ChebyshevInequalityOnWindows}
\noindent and the input needed at a given scale is the existence of level-compatible
(\texttt{EFamC}), tied families with windows:
\leanquote{Formalism.HC.ChebyshevInequalityAtWindowScale}
\noindent With $\varepsilon$ free, these give zero-freeness on $\operatorname{Re}s>\frac12$:
\leanquote{Formalism.HC.zeroFreeOn_of_tiedWin_cyc}
\noindent (In \lean{Formalism.HC.level_of_tiedWin} the proof picks $\varepsilon$ with
$(e_D+e_L+7)\varepsilon$ below half the distance from $\operatorname{Re}s$ to $\frac12$ and applies
the criterion with $\theta=\frac12+\frac12\min(\operatorname{Re}s-\frac12,\frac12)$
\stl{src/Formalism/SharpRestriction/RamificationFrobenius/TotientGaloisClasses/TopWin.lean}{68}.)

The first hypothesis holds for every level, with \texttt{Tied} the predicate
\lean{Formalism.MainNotions.ChainTied}:
\leanquote{UpperTheory.Comparison₂.level_gm_le_on_family_scales}
\noindent So, as written: for any level $\Lambda$ and $\varepsilon>0$, a family satisfying
\texttt{EFamC Λ} and \texttt{ChainTied} at all sufficiently large points of its window set has
$\Lambda.\mathrm{Gm}(x)\le x^{1/2+(e_D+e_L+7)\varepsilon}$ at all sufficiently large points of that
set. The difficulty is moved into producing level-compatible, chain-tied families. The predicate
\lean{Formalism.MainNotions.ChainTied}
\stl{src/Formalism/SharpRestriction/GlobalSchurBounds/TieGlobal.lean}{97} is a large compatibility
condition that ties the family's data at $x$ to: ranks and slopes $\mu$, $\mu_{\max}$ of lattices
(\texttt{TieLocal}), arithmetic data including the values \texttt{logDerivativeAtOne χ} of primitive
characters (\texttt{TieArith}), size data of an ``$L$-function chain''
(\texttt{LFunctionChainParameters.LChain.SizeData₂}) and error terms on windows (\texttt{ErrWindows}).
The slope vocabulary is that of Arakelov-theoretic slope inequalities, which remark item 13 names
(the Bost slope inequality) \sti~\stc{bost}. We have not unpacked this predicate.

\subsection{Where the families come from}

The remaining input is supplied for every $\psi$ by the following theorem, applied with
$\lambda=\frac32$:
\leanquote{UpperTheory.Top.zeroFreeAll_of_topG}
\noindent It has two hypotheses. The first is \lean{UpperTheory.Top.TopG}. For any $\lambda>1$, it
applies to a requirement \texttt{G} supplied by \texttt{HostExW λ G.req} whose class satisfies the
Galois-cohomology condition \lean{UpperTheory.Top.HypV}. It then supplies an eventually chain-tied
member of that class for each sufficiently small positive $\varepsilon$, with common exponents
$e_D,e_L$ (the class is ``connected''):
\leanquote{UpperTheory.Top.TopG}
\leanquote{UpperTheory.Top.FamilyClassConnected}
\noindent For the requirement \lean{UpperTheory.Top.GChi} $\psi$
\stl{src/Development/AutomorphicArithmetic/CharacterBoxes/Supply.lean}{42}, membership in the class
supplies the windows and level compatibility (\lean{UpperTheory.Top.eFamC_of_cls}); the requirement
combines an analytic condition at the cyclotomic level of $\psi$ with the Galois-cohomology bound
\lean{UpperTheory.Top.VVle} of \S\ref{sec:window}. The second input is
\texttt{HostExW λ (GChi ψ).req} for every primitive $\psi$. Both hypotheses are discharged in
\lean{UpperTheory.Top.zero_free_all_of_green_v_win}
\stl{src/Development/AutomorphicArithmetic/ThetaCorrespondence/ThetaPairLifts/ThirdPlaceSeries/TopPlug.lean}{12}:
\begin{itemize}
\item \lean{UpperTheory.Top.TopG} from Green's classification of cuspidal characters of finite
  general linear groups together with the theta-integral hypotheses (\S\ref{sec:thetaside}), through
  \lean{UpperTheory.Top.shape_closed_link_n_of_green}
  \stl{src/Development/AutomorphicArithmetic/ThetaCorrespondence/ThetaPairLifts/ThirdPlaceSeries/TopShape.lean}{33};
\item the supply from the window bound, through \lean{UpperTheory.Top.supplyChi}
  \stl{src/Development/AutomorphicArithmetic/CharacterBoxes/SupplyChi.lean}{8}:
  $\texttt{VWinAt}\ \psi\ 401\Rightarrow\texttt{HostExW}(3/2)\,(\texttt{GChi}\ \psi)\texttt{.req}$.
\end{itemize}

\paragraph{Summary of the analytic half.} For each primitive $\psi$, as written:
\[
\begin{aligned}
\underbrace{\text{\begin{tabular}{@{}c@{}}level-compatible, chain-tied\\families on windows\end{tabular}}}_{\text{TopG + supply}}
&\ \xRightarrow{\ \texttt{level\_gm\_le\_on\_family\_scales}\ }\
\Lambda.\mathrm{Gm}(x)\le x^{1/2+O(\varepsilon)}\text{ on windows}\\
&\ \xRightarrow{\ \text{Statement \ref{st:landau}}\ }\
L(s,\psi)\ne0\ (\operatorname{Re}s>\tfrac12).
\end{aligned}
\]
